\section{Sharpness and exact zeros}\label{sec:sharpness}
The parity restriction is sharp. If $d$ is even, pair the coordinates and
set $J(x_1,x_2,\ldots)=(-x_2,x_1,-x_4,x_3,\ldots)$. Each pair
contributes zero to $x\cdot Jx$, and $J$ preserves the Euclidean norm.
Thus there is an explicit unit-length tangent field on every
odd-dimensional unit sphere. At the other endpoint of the dimension
range, $S^0$ has zero tangent spaces and an exact zero at either point.

The uniform-modulus formulation also applies to fields that fail to be
Lipschitz. For example on $S^2$ put
\[
       F(x)=(1+\sqrt{|x_1|})\,(e_1-x_1 x).
\]
The second factor is tangent and has $\ell^1$ norm at most $1+d$.
Its change is bounded by $(1+d)\len{x-y}$, using $|x_1|\le1$ and
$\len x\le d$. The change of the square-root factor is at most
$\sqrt{\len{x-y}}$. Consequently for $\len{x-y}\le u\le1$ the total
change is at most $3(1+d)\sqrt{u}$. This gives, for example,
$\nu(h)=2(h+\clog(3(1+d)))$. Its values are sampled by bounded bisection.
At $(0,1,0)$ the first component, along
$(s,\sqrt{1-s^2},0)$ for small $s>0$, changes from $1$ to
$(1+\sqrt{s})(1-s^2)$, while the input distance is at most $2s$.
Thus no finite Lipschitz constant suffices.

\subsection{What a uniform exact-zero procedure would decide}
Even for $S^2$, selecting an exact zero of a uniformly given field
would have a logical consequence. Let $b(0),b(1),\ldots$ be binary
with at most one $1$. Form
\[
 u=(a,c,0),\qquad a=\sum_{s\ge0}2^{-s-1}b(2s),\qquad
 c=\sum_{s\ge0}2^{-s-1}b(2s+1),
 \qquad V(x)=u-(u\cdot x)x .
\]
The first $l$ terms give each coefficient to within $2^{-l}$.
Truncating at $l=h+\clog(2(1+d))+1$ is more than enough for
$\ell^1$ error $2^{-h}$ in $V$. All fields in this family have a
common continuity modulus, since $\len u\le1$ and
$\len{V(x)-V(y)}\le(1+d)\len{x-y}$. The field is smooth as a
function of $x$.

Suppose an exact zero $x$ is supplied by a fast name. Read a rational
$v$ within $1/8$ of $|x_1|$. If $v\ge1/2$, an odd occurrence of $1$
is impossible: it would make $u$ a positive multiple of $e_2$ and
force every zero to have $x_1=0$. Thus all odd bits are zero.
If $v<1/2$, an even occurrence would make $u$ a positive multiple
of $e_1$ and force $|x_1|=1$. Thus all even bits are zero. These are
constructive implications by testing each individual bit.
The resulting disjunction is the lesser limited principle of
omniscience (LLPO) on this input.

There is already a computability obstruction for a family of PR input
terms with a number parameter. Fix a machine enumeration $\varphi_i$.
Let $b_i(2s+j)$, $j=0,1$, indicate that machine $i$ on input $i$
first halts at step $s$ with output $j$. The bounded simulation is PR.
The reported parity side would give a total computable bit $c(i)$ equal
to $\varphi_i(i)$ whenever the latter is $0$ or $1$: report $0$ when
all odd bits are zero and $1$ in the other case. A program for
$i\mapsto1-c(i)$ has an index $h$; then $c(h)=1-c(h)$. This rules
out a uniform computable exact-zero selector for this family. Each one
of these particular fields individually has a computable pole. The
argument concerns uniform selection; the following example concerns a
single instance.

\subsection{A single primitive-recursively presented field whose poles have no primitive-recursive fast names}
\label{sec:no-pr-pole}

The following separation concerns the standard computability-theoretic
interpretation of iPRA. By a PR fast name of a point we mean a
primitive-recursive rational sequence with error at most $2^{-h}$ at
input $h$. We exhibit one field on $S^2$ with a PR sampler and a linear
modulus, for which neither of its two zeros has such a name.

Enumerate unary PR programs by natural numbers, treating an ill-formed
code as the constant-zero program. A Turing interpreter can evaluate any
program in this enumeration. For program $i$ write $f_i$ for its
function and put $b_i=1-\min(1,f_i(i))$. The binary sequence $(b_i)$
is computable. If it were PR, it would equal some $f_e$ and we would
have $b_e=1-b_e$. Define the real
\[
                     \alpha=\sum_{i=0}^{\infty}b_i4^{-i-1}
                                \in[0,1/3].
\]
To recover $b_i$ from a $2^{-(2i+5)}$-approximation of $\alpha$,
multiply that approximation by $4^{i+1}$ and round to the nearest
integer. The exact scaled real equals the integer coding the first
$i+1$ base-four digits plus a tail in $[0,1/3]$. The scaled error is
at most $1/8$, so rounding recovers that integer. Its residue modulo
four is $b_i$. Thus $\alpha$ has no PR fast name.

Fix a Turing machine that prints $b_0,b_1,\ldots$ successively.
After $s$ steps, let $j_s$ be the number of emitted bits and $l_s$
the base-four prefix sum of these bits. Define
\[
 I_s=[l_s,l_s+4^{-j_s}/3],\qquad
 g(t)=\sum_{s=0}^\infty 2^{-s-1}\operatorname{clamp}_{I_s}(t),
 \qquad \operatorname{clamp}_{[a,b]}(t)=\min(b,\max(a,t)).
\]
The endpoints of $I_s$ are PR in $s$: one simulates exactly $s$ steps
of a fixed machine, recording each output bit in its configuration.
One transition and iteration for a specified number of transitions are
PR. No bound on the waiting time for the next bit enters this operation.
The intervals are nested and contain $\alpha$. Their lengths tend to
zero because the machine eventually prints each bit. The sum defining
$g$ has values in $[0,1/3]$, admits a rational approximation on rational
inputs with error at most $2^{-M}/3$ from its first $M$ terms, and is
$1$-Lipschitz since each clamp has that property. Its unique fixed
point is $\alpha$. Indeed, for $t<\alpha$ every clamp is at least $t$,
and some $I_s$ has lower endpoint strictly above $t$. The corresponding
summand in $g(t)-t$ is strictly positive. The case $t>\alpha$ follows
from the upper endpoints, and $g(\alpha)=\alpha$ holds term by term.

On $P=\{p\in\mathbb R^3:\|p\|_1=1\}$ define
\[
 t(p)=\frac{p_1p_2}{\max(1/4,p_1^2)},\quad
 u(p)=(1,g(t(p)),0),\quad
 F(p)=u(p)-\frac{p\cdot u(p)}{p\cdot p}\,p .
\]
The denominator $p\cdot p$ is at least $1/3$. We give explicit
continuity and sampling bounds to verify the presentation.
For $\delta=\|p-q\|_1$, the numerator of $t$ changes by at most
$\delta$ and its denominator by at most $2\delta$. The denominators
are at least $1/4$, so $|t(p)-t(q)|\leq36\delta$ and
$\|u(p)-u(q)\|_1\leq36\delta$. Write $a=p\cdot p$ and
$b=p\cdot u(p)$. Since $\|u(p)\|_\infty=1$, we have $|b|\leq1$,
the change of $b$ is at most $37\delta$, and the change of $a$ is
at most $2\delta$. Hence $|b/a|\leq3$ and the change of $b/a$ is
at most $(3\cdot37+9\cdot2)\delta=129\delta$. Substituting into
$F$ gives $\|F(p)-F(q)\|_1\leq168\delta$.

For $A(p,h)$ evaluate the displayed rational formula at rational
$p$, using the first $M=h+1$ terms of $g$. The projection
$v\mapsto v-(p\cdot v)p/(p\cdot p)$ has $\ell^1$ operator norm
at most $4$, so the sampler error is at most $4\cdot2^{-(h+1)}/3
\leq 2^{-h}$. Its dot product with $p$ vanishes exactly. Two such
sample errors give the coherence bound, and the Lipschitz constant
$168\leq 2^8$ gives the sampled continuity bound with $\mu(h)=h+8$.
The arithmetic verification can also be made directly on truncations:
their ranges are in $[0,1/3]$, each has Lipschitz constant at most one,
and the difference of a longer and a shorter truncation is bounded by
the omitted geometric tail. These finite checks require no proof in
iPRA that the printing machine eventually emits every digit. Its
eventual progress is used in the external argument identifying the poles.

If $F(p)=0$ then $p$ and $u(p)$ are parallel. Since the first
coordinate of $u(p)$ is $1$, we have $p_1\ne0$,
$p_3=0$, $p_2/p_1=g(t(p))$ and $|p_1|=1/(1+g(t(p)))\geq3/4$.
The denominator in $t(p)$ is then $p_1^2$, proving
$t(p)=p_2/p_1=g(t(p))$. Consequently $t(p)=\alpha$.
Conversely the two poles
$\pm(1,\alpha,0)/(1+\alpha)$ satisfy $F(p)=0$.
A PR fast name of either pole gives one of $t(p)$ by the displayed
rational formula: on the coordinate cube $[-1,1]^3$ the same
$36$-Lipschitz estimate applies, and rational approximations can first
be clipped to the cube. This contradicts the decoding argument.

On the round sphere use the tangent field $x\mapsto F(x/\|x\|_1)$.
The radial map is $2$-Lipschitz in $\ell^1$ there, since
$\|x\|_1\geq1$. It carries a PR-named zero on the round sphere to a
PR-named zero on $P$. The two poles do have ordinary computable names,
since the bit-printing program eventually evaluates every digit.
The separation is precisely between PR finite-precision witnesses
and a PR fast name for an exact zero of this one field.

